\section{实验结果与可操作结论}

\subsection{结果解读：趋势比绝对数值更重要}
课程展示了在多项 benchmark 上的提升趋势，尤其在更难任务和在线设置下，加入规划与更强训练数据治理后效果更稳。即便具体分数随配置变化，趋势仍非常一致。

\begin{figure}[H]
\centering
\includegraphics[width=0.88\textwidth]{frame-08-results.jpg}
\caption{视频约 01:24:55：结果页强调分阶段训练与 planner 对复杂任务成功率的拉升}
\end{figure}

\begin{importantbox}{可复用的三条工程经验}
\begin{enumerate}
    \item 先把评估与环境打稳，再谈模型迭代速度；
    \item 数据治理优先级不低于模型升级；
    \item 对复杂任务，规划信号通常是 ``必要条件'' 而不是 ``锦上添花''。
\end{enumerate}
\end{importantbox}

\subsection{落地指标建议}
\begin{table}[H]
\centering
\begin{tabular}{p{3.2cm}p{5.0cm}p{2.8cm}}
\toprule
\textbf{指标类别} & \textbf{建议指标} & \textbf{使用场景} \\
\midrule
任务结果 & 完成率、部分完成率、失败类型分布 & 评估产品可用性 \\
轨迹质量 & 平均步数、重复动作率、回退率 & 诊断规划质量 \\
系统稳定性 & 重跑一致性、异常中断恢复率 & 评估生产可靠性 \\
数据闭环 & 失败样本回流比例、修复后提升幅度 & 衡量迭代效率 \\
\bottomrule
\end{tabular}
\caption{从课程结论抽象出的 Agent 工程指标框架}
\end{table}

\begin{warningbox}{不要只看 ``总完成率''}
单一指标会掩盖问题。比如完成率上升可能来自简单任务占比增加，而不是复杂任务能力提升。课程建议至少按任务难度、任务类型、环境类型做分层统计。
\end{warningbox}

\subsection{本章小结}
本讲结果最有价值的部分是方法论：环境、评估、数据、模型四者必须同时推进。任何一环短板都会让系统在真实任务上失稳。

